\begin{table}[htbp]
\centering
\small
\caption{§B-1 行政区划锚定 --- 三线建设县样本框架}
\label{tab:anchoring}
\begin{tabular}{l c c l}
\toprule
 & 锚定年份 & 县数 & 说明 \\
\midrule
\textbf{主样本 (Main)} & 1985 & \textbf{216} & 含 1997 重庆剥离后仍属川黔的县 (Hard Anchor) \
扩展 1 (含重庆) & 1985 & \textbf{249} & 含已剥离的 30+ 重庆县 (扣重) \
扩展 2 (2010 Soft Tracking) & 2010 & $\approx 233$ & 含撤县设区, 应用 spatio-temporal 重分配 \
\bottomrule
\end{tabular}
\begin{quote}\footnotesize
\textbf{主要区划变动 (1978-2017):}
\begin{itemize}
\item 1983: 涪陵/泸州/绵阳/乐山/内江升地级市 (县级边界基本未变, 仅行政归属调整)
\item 1989-1996: 重庆升级计划单列市, 川东 12 县改重庆代管 (仍属四川省口径)
\item 1997: 重庆直辖, 涪陵/万县/黔江/部分重庆区 $\to$ 重庆 (本研究 Hard Anchor 视为'含重庆')
\item 1997-2000: 四川眉山/资阳/巴中/广安 撤地设市 (县级边界不变)
\item 2000-2014: 贵州多轮撤地设市 (遵义 1997, 安顺 2000, 毕节/铜仁 2011)
\item 2010-2017: 撤县设区密集期 (如贵阳清镇/遵义播州/六盘水水城改区)
\end{itemize}
\textit{策略: 主样本用 1985 Hard Anchor (216 县), 检验样本 1997/2010 Soft Tracking.}
\end{quote}
\end{table}